\input format.tex

%\usepackage{arydshln}
\usepackage{graphicx}
\graphicspath{{cores/}}

\begin{document}

%%%
\setlength{\arrayrulewidth}{0.5pt}
\fontsize{9.3pt}{17pt}\selectfont
\color{gray2}

%%%% protein
\vspace*{0mm}
\begin{center}
{\bf\sanhao 物质代谢~-~蛋白质的消化与吸收}
\end{center}


\medskip
\noindent
{\bf\xiaosihao Tom先生：}

\bigskip

蛋白质是组成人体细胞、组织的重要成分，是生命活动的主要承担者，没有蛋白质就没有生命。除了与结构和性状有关，蛋白质还参与基因表达的调节、各种生物化学反应、神经传递等生命过程。

您肠道内参与蛋白质总量、谷胱甘肽代谢和牛磺酸代谢的菌群指标均处于正常水平，请您继续保持良好的饮食及生活习惯。

\fontsize{9.3pt}{9.8pt}\selectfont
\bigskip
{
\lantxh
%\begin{tctabularx}{tabularx={m{3.75cm}m{4.9cm}<{\centering}m{1.22cm}<{\centering}m{1.8cm}<{\centering}m{2cm}<{\centering}}}
\begin{tctabularx}{tabularx={m{4.8cm}m{5.2cm}<{\centering}m{4.55cm}<{\centering}}}
%\cellcolor{kelly} &
%\cellcolor{kelly} &
%\cellcolor{kelly} &
%\cellcolor{kelly} &
%\cellcolor{kelly}
&&
\\[-6pt]
\cellcolor{kelly} \raisebox{2.614pt}{\color{black70}\bf 测试项目} &
\cellcolor{kelly} \raisebox{2.614pt}{\color{black70}\bf 分布} &
%\cellcolor{kelly} \raisebox{2.614pt}{\color{black70}\bf 检测值} &
%\cellcolor{kelly} \raisebox{2.614pt}{\color{black70}\bf 参考范围} &
\cellcolor{kelly} \raisebox{2.614pt}{\color{black70}\bf 检测结果}
%{\color{black70}\bf 测试项目} & {\color{black70}\bf 分布} & {\color{black70}\bf 检测值} &{\color{black70}\bf 参考范围} & {\color{black70}\bf 检测结果}
\end{tctabularx}

\vspace*{-4.25mm}
\fontsize{8.8pt}{11pt}\selectfont
%\begin{longtable}{m{3.75cm}m{5cm}<{\centering}m{0cm}@{}m{1.22cm}<{\centering}m{1.8cm}<{\centering}m{2cm}<{\centering}}
\begin{longtable}{m{4.8cm}m{5.2cm}<{\centering}m{0cm}@{}m{4.61cm}<{\centering}}
\hline
%\parbox[c]{\hsize}{\vskip7pt 蛋白质总量\\蛋白质是一切生命活动的物质基础 \vskip7pt} & \parbox[c]{\hsize}{\vskip7pt\centerline{\raisebox{-1.5ex}{\includegraphics[width=5cm,height=0.3cm]{zxdanbaizhizongliang_dis.png}}}\vskip7pt}  &\hspace*{-4.42347098cm}\raisebox{-0.45ex}{\includegraphics{triangle.pdf}} & \begin{minipage}{4.60cm}\begin{center}{正常 }\end{center} \end{minipage} \\
\parbox[c]{\hsize}{\vskip7pt 蛋白质总量\\蛋白质是一切生命活动的物质基础 \vskip7pt} & \parbox[c]{\hsize}{\vskip7pt\centerline{\raisebox{-1.5ex}{\includegraphics[width=5cm,height=0.3cm]{zxdanbaizhizongliang_dis.png}}}\vskip7pt}  &\hspace*{-4.42347098cm}\raisebox{-0.45ex}{\includegraphics{triangle.pdf}} & \begin{minipage}{4.60cm}\begin{center}{正常\\ \bahao 蛋白质均衡 }\end{center} \end{minipage} \\
\hline
%\parbox[c]{\hsize}{\vskip7pt 谷胱甘肽代谢\\抗衰老，解毒 \vskip7pt} & \parbox[c]{\hsize}{\vskip7pt\centerline{\raisebox{-1.5ex}{\includegraphics[width=5cm,height=0.3cm]{zxguguanggantaidaixie_dis.png}}}\vskip7pt}  &\hspace*{-2.163835322cm}\raisebox{-0.45ex}{\includegraphics{triangle.pdf}} & \begin{minipage}{4.60cm}\begin{center}{正常 }\end{center} \end{minipage} \\
\parbox[c]{\hsize}{\vskip7pt 谷胱甘肽代谢\\抗衰老，解毒 \vskip7pt} & \parbox[c]{\hsize}{\vskip7pt\centerline{\raisebox{-1.5ex}{\includegraphics[width=5cm,height=0.3cm]{zxguguanggantaidaixie_dis.png}}}\vskip7pt}  &\hspace*{-2.163835322cm}\raisebox{-0.45ex}{\includegraphics{triangle.pdf}} & \begin{minipage}{4.60cm}\begin{center}{正常 }\end{center} \end{minipage} \\
\hline
%\parbox[c]{\hsize}{\vskip7pt 牛磺酸代谢\\维持心血管、骨骼肌、视网膜和中枢神经系统正常功能 \vskip7pt} & \parbox[c]{\hsize}{\vskip7pt\centerline{\raisebox{-1.5ex}{\includegraphics[width=5cm,height=0.3cm]{zxniuhuangsuandaixie_dis.png}}}\vskip7pt}  &\hspace*{-3.414116612cm}\raisebox{-0.45ex}{\includegraphics{triangle.pdf}} & \begin{minipage}{4.60cm}\begin{center}{正常 }\end{center} \end{minipage} \\
\parbox[c]{\hsize}{\vskip7pt 牛磺酸代谢\\维持心血管、骨骼肌、视网膜和中枢神经系统正常功能 \vskip7pt} & \parbox[c]{\hsize}{\vskip7pt\centerline{\raisebox{-1.5ex}{\includegraphics[width=5cm,height=0.3cm]{zxniuhuangsuandaixie_dis.png}}}\vskip7pt}  &\hspace*{-3.414116612cm}\raisebox{-0.45ex}{\includegraphics{triangle.pdf}} & \begin{minipage}{4.60cm}\begin{center}{正常 }\end{center} \end{minipage} \\
\hline
\end{longtable}

\noindent
\tikz\draw[green2,fill=green2](0,0) circle[radius=4pt]; 区域表示正常人群的一般水平 \tikz\draw[darkblue,fill=darkblue](0,0) circle[radius=4pt]; 区域表示仍在正常水平，但较一般水平偏高或偏低 %\tikz\draw[red,fill=red](0,0) circle[radius=4pt]; 而红色表示为有风险
}

\bigskip
{\centering {\includegraphics{jieshi.pdf}}}

%%% fat
\newpage
\setlength{\arrayrulewidth}{0.5pt}
\fontsize{9.3pt}{17pt}\selectfont
\color{gray2}

\vspace*{0mm}
\begin{center}
{\bf\sanhao 物质代谢~-~脂类的消化与吸收}
\end{center}

\medskip
\noindent
{\bf\xiaosihao Tom先生：}

\bigskip

脂类是人体主要储能物质，同时起到保护脏器、保持体温的作用。除此之外，脂类还参与细胞膜的构成、信号传递、形成激素等生命过程。

您肠道内参与花生四烯酸代谢的菌群指标优于参考人群，有利于脂质蛋白的代谢、防治心血管疾病；您肠道内参与甘油三酯代谢、鞘脂代谢、亚油酸代谢的菌群指标存在异常，可能增加动脉硬化、心脑血管疾病的风险，不利于细胞的生长与凋亡，需引起注意，建议持续监测肠道健康。

\fontsize{9.3pt}{9.8pt}\selectfont
\bigskip
{
\lantxh
%\begin{tctabularx}{tabularx={m{3.75cm}m{4.9cm}<{\centering}m{1.22cm}<{\centering}m{1.8cm}<{\centering}m{2cm}<{\centering}}}
\begin{tctabularx}{tabularx={m{4.8cm}m{5.2cm}<{\centering}m{4.55cm}<{\centering}}}
%\cellcolor{kelly} &
%\cellcolor{kelly} &
%\cellcolor{kelly} &
%\cellcolor{kelly} &
%\cellcolor{kelly}
&&
\\[-6pt]
\cellcolor{kelly} \raisebox{2.614pt}{\color{black70}\bf 测试项目} &
\cellcolor{kelly} \raisebox{2.614pt}{\color{black70}\bf 分布} &
%\cellcolor{kelly} \raisebox{2.614pt}{\color{black70}\bf 检测值} &
%\cellcolor{kelly} \raisebox{2.614pt}{\color{black70}\bf 参考范围} &
\cellcolor{kelly} \raisebox{2.614pt}{\color{black70}\bf 检测结果}
%{\color{black70}\bf 测试项目} & {\color{black70}\bf 分布} & {\color{black70}\bf 检测值} &{\color{black70}\bf 参考范围} & {\color{black70}\bf 检测结果}
\end{tctabularx}

\vspace*{-4.25mm}
\fontsize{8.8pt}{11pt}\selectfont
%\begin{longtable}{m{3.75cm}m{5cm}<{\centering}m{0cm}@{}m{1.22cm}<{\centering}m{1.8cm}<{\centering}m{2cm}<{\centering}}
\begin{longtable}{m{4.8cm}m{5.2cm}<{\centering}m{0cm}@{}m{4.61cm}<{\centering}}
\hline
%\parbox[c]{\hsize}{\vskip7pt 脂肪总量\\保护脏器，提供必需脂肪酸 \vskip7pt} & \parbox[c]{\hsize}{\vskip7pt\centerline{\raisebox{-1.5ex}{\includegraphics[width=5cm,height=0.3cm]{zxzhifangzongliang_dis.png}}}\vskip7pt}  &\hspace*{-3.250211366cm}\raisebox{-0.45ex}{\includegraphics{triangle.pdf}} & \begin{minipage}{4.60cm}\begin{center}{正常 }\end{center} \end{minipage} \\
\parbox[c]{\hsize}{\vskip7pt 脂肪总量\\保护脏器，提供必需脂肪酸 \vskip7pt} & \parbox[c]{\hsize}{\vskip7pt\centerline{\raisebox{-1.5ex}{\includegraphics[width=5cm,height=0.3cm]{zxzhifangzongliang_dis.png}}}\vskip7pt}  &\hspace*{-3.250211366cm}\raisebox{-0.45ex}{\includegraphics{triangle.pdf}} & \begin{minipage}{4.60cm}\begin{center}{正常\\ \bahao 脂肪代谢均衡 }\end{center} \end{minipage} \\
\hline
%\parbox[c]{\hsize}{\vskip7pt 脂肪酸代谢\\降脂、抗动脉硬化 \vskip7pt} & \parbox[c]{\hsize}{\vskip7pt\centerline{\raisebox{-1.5ex}{\includegraphics[width=5cm,height=0.3cm]{zxzhifangsuandaixie_dis.png}}}\vskip7pt}  &\hspace*{-3.434925542cm}\raisebox{-0.45ex}{\includegraphics{triangle.pdf}} & \begin{minipage}{4.60cm}\begin{center}{正常 }\end{center} \end{minipage} \\
\parbox[c]{\hsize}{\vskip7pt 脂肪酸代谢\\降脂、抗动脉硬化 \vskip7pt} & \parbox[c]{\hsize}{\vskip7pt\centerline{\raisebox{-1.5ex}{\includegraphics[width=5cm,height=0.3cm]{zxzhifangsuandaixie_dis.png}}}\vskip7pt}  &\hspace*{-3.434925542cm}\raisebox{-0.45ex}{\includegraphics{triangle.pdf}} & \begin{minipage}{4.60cm}\begin{center}{正常 }\end{center} \end{minipage} \\
\hline
%\parbox[c]{\hsize}{\vskip7pt 甘油三酯代谢\\为机体提供能量，与动脉硬化有关 \vskip7pt} & \parbox[c]{\hsize}{\vskip7pt\centerline{\raisebox{-1.5ex}{\includegraphics[width=5cm,height=0.3cm]{zxganyousanzhidaixie_dis.png}}}\vskip7pt}  &\hspace*{-5.5236753386cm}\raisebox{-0.45ex}{\includegraphics{triangle.pdf}} & \begin{minipage}{4.60cm}\begin{center}{偏低 \\ \bahao 供能减少，可能增加动脉硬化风险 }\end{center} \end{minipage} \\
\parbox[c]{\hsize}{\vskip7pt 甘油三酯代谢\\为机体提供能量，与动脉硬化有关 \vskip7pt} & \parbox[c]{\hsize}{\vskip7pt\centerline{\raisebox{-1.5ex}{\includegraphics[width=5cm,height=0.3cm]{zxganyousanzhidaixie_dis.png}}}\vskip7pt}  &\hspace*{-5.5236753386cm}\raisebox{-0.45ex}{\includegraphics{triangle.pdf}} & \begin{minipage}{4.60cm}\begin{center}{偏低\\ \bahao 供能减少，可能增加动脉硬化风险 }\end{center} \end{minipage} \\
\hline
%\parbox[c]{\hsize}{\vskip7pt 甘油磷脂代谢\\降脂、防治心脑血管病，增强记忆力 \vskip7pt} & \parbox[c]{\hsize}{\vskip7pt\centerline{\raisebox{-1.5ex}{\includegraphics[width=5cm,height=0.3cm]{zxganyoulinzhidaixie_dis.png}}}\vskip7pt}  &\hspace*{-4.379127068cm}\raisebox{-0.45ex}{\includegraphics{triangle.pdf}} & \begin{minipage}{4.60cm}\begin{center}{正常 }\end{center} \end{minipage} \\
\parbox[c]{\hsize}{\vskip7pt 甘油磷脂代谢\\降脂、防治心脑血管病，增强记忆力 \vskip7pt} & \parbox[c]{\hsize}{\vskip7pt\centerline{\raisebox{-1.5ex}{\includegraphics[width=5cm,height=0.3cm]{zxganyoulinzhidaixie_dis.png}}}\vskip7pt}  &\hspace*{-4.379127068cm}\raisebox{-0.45ex}{\includegraphics{triangle.pdf}} & \begin{minipage}{4.60cm}\begin{center}{正常 }\end{center} \end{minipage} \\
\hline
%\parbox[c]{\hsize}{\vskip7pt 鞘脂代谢\\参与调节细胞的生长与凋亡 \vskip7pt} & \parbox[c]{\hsize}{\vskip7pt\centerline{\raisebox{-1.5ex}{\includegraphics[width=5cm,height=0.3cm]{zxqiaozhidaixie_dis.png}}}\vskip7pt}  &\hspace*{-5.07532808cm}\raisebox{-0.45ex}{\includegraphics{triangle.pdf}} & \begin{minipage}{4.60cm}\begin{center}{偏低 \\ \bahao 不利于细胞的生长与凋亡 }\end{center} \end{minipage} \\
\parbox[c]{\hsize}{\vskip7pt 鞘脂代谢\\参与调节细胞的生长与凋亡 \vskip7pt} & \parbox[c]{\hsize}{\vskip7pt\centerline{\raisebox{-1.5ex}{\includegraphics[width=5cm,height=0.3cm]{zxqiaozhidaixie_dis.png}}}\vskip7pt}  &\hspace*{-5.07532808cm}\raisebox{-0.45ex}{\includegraphics{triangle.pdf}} & \begin{minipage}{4.60cm}\begin{center}{偏低\\ \bahao 不利于细胞的生长与凋亡 }\end{center} \end{minipage} \\
\hline
%\parbox[c]{\hsize}{\vskip7pt 亚油酸代谢\\预防心脑血管疾病，抗癌抗过敏，健脑，保护视力 \vskip7pt} & \parbox[c]{\hsize}{\vskip7pt\centerline{\raisebox{-1.5ex}{\includegraphics[width=5cm,height=0.3cm]{zxyayousuandaixie_dis.png}}}\vskip7pt}  &\hspace*{-5.3819794982cm}\raisebox{-0.45ex}{\includegraphics{triangle.pdf}} & \begin{minipage}{4.60cm}\begin{center}{偏低 \\ \bahao 可能增加心脑血管疾病及过敏风险 }\end{center} \end{minipage} \\
\parbox[c]{\hsize}{\vskip7pt 亚油酸代谢\\预防心脑血管疾病，抗癌抗过敏，健脑，保护视力 \vskip7pt} & \parbox[c]{\hsize}{\vskip7pt\centerline{\raisebox{-1.5ex}{\includegraphics[width=5cm,height=0.3cm]{zxyayousuandaixie_dis.png}}}\vskip7pt}  &\hspace*{-5.3819794982cm}\raisebox{-0.45ex}{\includegraphics{triangle.pdf}} & \begin{minipage}{4.60cm}\begin{center}{偏低\\ \bahao 可能增加心脑血管疾病及过敏风险 }\end{center} \end{minipage} \\
\hline
%\parbox[c]{\hsize}{\vskip7pt 花生四烯酸代谢\\调节脂质蛋白的代谢，有助于预防心脑血管疾病 \vskip7pt} & \parbox[c]{\hsize}{\vskip7pt\centerline{\raisebox{-1.5ex}{\includegraphics[width=5cm,height=0.3cm]{zxhuashengsixisuandaixie_dis.png}}}\vskip7pt}  &\hspace*{-0.737811326cm}\raisebox{-0.45ex}{\includegraphics{triangle.pdf}} & \begin{minipage}{4.60cm}\begin{center}{高\\ \bahao 有利于脂质蛋白的代谢、防治心血管疾病 }\end{center} \end{minipage} \\
\parbox[c]{\hsize}{\vskip7pt 花生四烯酸代谢\\调节脂质蛋白的代谢，有助于预防心脑血管疾病 \vskip7pt} & \parbox[c]{\hsize}{\vskip7pt\centerline{\raisebox{-1.5ex}{\includegraphics[width=5cm,height=0.3cm]{zxhuashengsixisuandaixie_dis.png}}}\vskip7pt}  &\hspace*{-0.737811326cm}\raisebox{-0.45ex}{\includegraphics{triangle.pdf}} & \begin{minipage}{4.60cm}\begin{center}{高\\ \bahao 有利于脂质蛋白的代谢、防治心血管疾病 }\end{center} \end{minipage} \\
\hline
\end{longtable}

\noindent
\tikz\draw[green2,fill=green2](0,0) circle[radius=4pt]; 区域表示正常人群的一般水平 \tikz\draw[darkblue,fill=darkblue](0,0) circle[radius=4pt]; 区域表示仍在正常水平，但较一般水平偏高或偏低 %\tikz\draw[red,fill=red](0,0) circle[radius=4pt]; 而红色表示为有风险
}

\bigskip
{\centering {\includegraphics{jieshi.pdf}}}

%%% carbon
\newpage

\setlength{\arrayrulewidth}{0.5pt}
\fontsize{9.3pt}{17pt}\selectfont
\color{gray2}

\vspace*{0mm}
\begin{center}
{\bf\sanhao 物质代谢~-~碳水化合物的消化与吸收}
\end{center}

\medskip
\noindent
{\bf\xiaosihao Tom先生：}

\bigskip

碳水化合物是人体直接供能物质，维持机体的正常活动。除此之外，碳水化合物还参与维持大脑功能、调节脂肪代谢等重要生命过程。

您肠道内参与碳水化合物总量、半乳糖代谢、果糖和甘露糖代谢等的菌群指标存在异常
，不利于抑制病原菌生长，需引起注意，建议持续监测肠道健康。

\fontsize{9.3pt}{9.8pt}\selectfont
\bigskip
{
\lantxh
%\begin{tctabularx}{tabularx={m{3.75cm}m{4.9cm}<{\centering}m{1.22cm}<{\centering}m{1.8cm}<{\centering}m{2cm}<{\centering}}}
\begin{tctabularx}{tabularx={m{4.8cm}m{5.2cm}<{\centering}m{4.55cm}<{\centering}}}
%\cellcolor{kelly} &
%\cellcolor{kelly} &
%\cellcolor{kelly} &
%\cellcolor{kelly} &
%\cellcolor{kelly}
&&
\\[-6pt]
\cellcolor{kelly} \raisebox{2.614pt}{\color{black70}\bf 测试项目} &
\cellcolor{kelly} \raisebox{2.614pt}{\color{black70}\bf 分布} &
%\cellcolor{kelly} \raisebox{2.614pt}{\color{black70}\bf 检测值} &
%\cellcolor{kelly} \raisebox{2.614pt}{\color{black70}\bf 参考范围} &
\cellcolor{kelly} \raisebox{2.614pt}{\color{black70}\bf 检测结果}
%{\color{black70}\bf 测试项目} & {\color{black70}\bf 分布} & {\color{black70}\bf 检测值} &{\color{black70}\bf 参考范围} & {\color{black70}\bf 检测结果}
\end{tctabularx}

\vspace*{-4.25mm}
\fontsize{8.8pt}{11pt}\selectfont
%\begin{longtable}{m{3.75cm}m{5cm}<{\centering}m{0cm}@{}m{1.22cm}<{\centering}m{1.8cm}<{\centering}m{2cm}<{\centering}}
\begin{longtable}{m{4.8cm}m{5.2cm}<{\centering}m{0cm}@{}m{4.61cm}<{\centering}}
\hline
%\parbox[c]{\hsize}{\vskip7pt 碳水化合物总量\\主要的供能物质 \vskip7pt} & \parbox[c]{\hsize}{\vskip7pt\centerline{\raisebox{-1.5ex}{\includegraphics[width=5cm,height=0.3cm]{zxtanshuihuahewuzongliang_dis.png}}}\vskip7pt}  &\hspace*{-5.63497391636cm}\raisebox{-0.45ex}{\includegraphics{triangle.pdf}} & \begin{minipage}{4.60cm}\begin{center}{低 \\ \bahao 不利于维持肠道稳态和菌群平衡 }\end{center} \end{minipage} \\
\parbox[c]{\hsize}{\vskip7pt 碳水化合物总量\\主要的供能物质 \vskip7pt} & \parbox[c]{\hsize}{\vskip7pt\centerline{\raisebox{-1.5ex}{\includegraphics[width=5cm,height=0.3cm]{zxtanshuihuahewuzongliang_dis.png}}}\vskip7pt}  &\hspace*{-5.63497391636cm}\raisebox{-0.45ex}{\includegraphics{triangle.pdf}} & \begin{minipage}{4.60cm}\begin{center}{低\\ \bahao 不利于维持肠道稳态和菌群平衡 }\end{center} \end{minipage} \\
\hline
%\parbox[c]{\hsize}{\vskip7pt 半乳糖代谢\\生成短链脂肪酸和乳酸，有利于肠道有益菌的生长 \vskip7pt} & \parbox[c]{\hsize}{\vskip7pt\centerline{\raisebox{-1.5ex}{\includegraphics[width=5cm,height=0.3cm]{zxbanrutangdaixie_dis.png}}}\vskip7pt}  &\hspace*{-5.562588959cm}\raisebox{-0.45ex}{\includegraphics{triangle.pdf}} & \begin{minipage}{4.60cm}\begin{center}{低 \\ \bahao 不利于半乳糖的消化及肠道有益菌的生长 }\end{center} \end{minipage} \\
\parbox[c]{\hsize}{\vskip7pt 半乳糖代谢\\生成短链脂肪酸和乳酸，有利于肠道有益菌的生长 \vskip7pt} & \parbox[c]{\hsize}{\vskip7pt\centerline{\raisebox{-1.5ex}{\includegraphics[width=5cm,height=0.3cm]{zxbanrutangdaixie_dis.png}}}\vskip7pt}  &\hspace*{-5.562588959cm}\raisebox{-0.45ex}{\includegraphics{triangle.pdf}} & \begin{minipage}{4.60cm}\begin{center}{低\\ \bahao 不利于半乳糖的消化及肠道有益菌的生长 }\end{center} \end{minipage} \\
\hline
%\parbox[c]{\hsize}{\vskip7pt 果糖和甘露糖代谢\\生成有机酸，抑制病原菌生长 \vskip7pt} & \parbox[c]{\hsize}{\vskip7pt\centerline{\raisebox{-1.5ex}{\includegraphics[width=5cm,height=0.3cm]{zxguotangheganlutangdaixie_dis.png}}}\vskip7pt}  &\hspace*{-5.5666672802cm}\raisebox{-0.45ex}{\includegraphics{triangle.pdf}} & \begin{minipage}{4.60cm}\begin{center}{低 \\ \bahao 不利于有机酸的产生，对病原菌的抑制作用减弱 }\end{center} \end{minipage} \\
\parbox[c]{\hsize}{\vskip7pt 果糖和甘露糖代谢\\生成有机酸，抑制病原菌生长 \vskip7pt} & \parbox[c]{\hsize}{\vskip7pt\centerline{\raisebox{-1.5ex}{\includegraphics[width=5cm,height=0.3cm]{zxguotangheganlutangdaixie_dis.png}}}\vskip7pt}  &\hspace*{-5.5666672802cm}\raisebox{-0.45ex}{\includegraphics{triangle.pdf}} & \begin{minipage}{4.60cm}\begin{center}{低\\ \bahao 不利于有机酸的产生，对病原菌的抑制作用减弱 }\end{center} \end{minipage} \\
\hline
%\parbox[c]{\hsize}{\vskip7pt 蔗糖和淀粉代谢\\产生短链脂肪酸，保护肠粘膜，供应能量 \vskip7pt} & \parbox[c]{\hsize}{\vskip7pt\centerline{\raisebox{-1.5ex}{\includegraphics[width=5cm,height=0.3cm]{zxzhetanghedianfendaixie_dis.png}}}\vskip7pt}  &\hspace*{-5.5305340436cm}\raisebox{-0.45ex}{\includegraphics{triangle.pdf}} & \begin{minipage}{4.60cm}\begin{center}{偏低 \\ \bahao 不利于保护肠粘膜 }\end{center} \end{minipage} \\
\parbox[c]{\hsize}{\vskip7pt 蔗糖和淀粉代谢\\产生短链脂肪酸，保护肠粘膜，供应能量 \vskip7pt} & \parbox[c]{\hsize}{\vskip7pt\centerline{\raisebox{-1.5ex}{\includegraphics[width=5cm,height=0.3cm]{zxzhetanghedianfendaixie_dis.png}}}\vskip7pt}  &\hspace*{-5.5305340436cm}\raisebox{-0.45ex}{\includegraphics{triangle.pdf}} & \begin{minipage}{4.60cm}\begin{center}{偏低\\ \bahao 不利于保护肠粘膜 }\end{center} \end{minipage} \\
\hline
%\parbox[c]{\hsize}{\vskip7pt 纤维素代谢\\分解纤维素产生短链脂肪酸等有益物质，促进肠道菌生长 \vskip7pt} & \parbox[c]{\hsize}{\vskip7pt\centerline{\raisebox{-1.5ex}{\includegraphics[width=5cm,height=0.3cm]{zxxianweisudaixie_dis.png}}}\vskip7pt}  &\hspace*{-1.804452128cm}\raisebox{-0.45ex}{\includegraphics{triangle.pdf}} & \begin{minipage}{4.60cm}\begin{center}{正常 }\end{center} \end{minipage} \\
\parbox[c]{\hsize}{\vskip7pt 纤维素代谢\\分解纤维素产生短链脂肪酸等有益物质，促进肠道菌生长 \vskip7pt} & \parbox[c]{\hsize}{\vskip7pt\centerline{\raisebox{-1.5ex}{\includegraphics[width=5cm,height=0.3cm]{zxxianweisudaixie_dis.png}}}\vskip7pt}  &\hspace*{-1.804452128cm}\raisebox{-0.45ex}{\includegraphics{triangle.pdf}} & \begin{minipage}{4.60cm}\begin{center}{正常 }\end{center} \end{minipage} \\
\hline
\end{longtable}

\noindent
\tikz\draw[green2,fill=green2](0,0) circle[radius=4pt]; 区域表示正常人群的一般水平 \tikz\draw[darkblue,fill=darkblue](0,0) circle[radius=4pt]; 区域表示仍在正常水平，但较一般水平偏高或偏低 %\tikz\draw[red,fill=red](0,0) circle[radius=4pt]; 而红色表示为有风险
}

\bigskip
{\centering {\includegraphics{jieshi.pdf}}}

\end{document}
